\documentclass[11pt,a4paper]{article}

% Page layout and mathematical tools
\usepackage[margin=1in]{geometry}
\usepackage{amsmath,amssymb,amsthm}
\usepackage[T1]{fontenc}
\usepackage{lmodern}
\usepackage{hyperref}

\title{Homological Algebra Lecture Notes}
\author{Ryland Gross}
\date{\today}

% Number results within each section. These environments share one counter.
\theoremstyle{plain}
\newtheorem{theorem}{Theorem}[section]
\newtheorem{proposition}[theorem]{Proposition}
\newtheorem{lemma}[theorem]{Lemma}
\newtheorem{corollary}[theorem]{Corollary}

\theoremstyle{definition}
\newtheorem{definition}[theorem]{Definition}
\newtheorem{example}[theorem]{Example}

\theoremstyle{remark}
\newtheorem{remark}[theorem]{Remark}

\begin{document}

\maketitle

\section{Lecture 1 (October 1st 2026)}

\noindent
\textbf{Preview:} We basically spend some time doing category theory, and the theory of functors with some attention paid to $R$-Modules. There is some additional mastery material that we need to cover as well.

\subsection*{Basics of Category Theory}

\begin{definition}
    A \emph{category} $\mathcal{C}$ consists of 
    \begin{enumerate}
        \item a collection of `objects' called $\text{ob}(\mathcal{C})$
        \item $\forall A, B \in \text{ob}(\mathcal{C})$ we have a set of ``morphisms'' $\text{Hom}_{\mathcal{C}}(A,B)$
        \item An identity (or a unit) for each object $\text{id}_A \in \text{Hom}_{\mathcal{C}}(A,A)$
        \item A composition rule for $A,B,C$ which is: $\text{Hom}_{\mathcal{C}}(B,C) \times \text{Hom}_{\mathcal{C}}(A,B) \to \text{Hom}_{\mathcal{C}}(A,C)$ It is simple to think of $g \circ f \in \text{Hom}_{\mathcal{C}}(A,C)$ where $g \in \text{Hom}_{\mathcal{C}}(B,C)$ and $f \in \text{Hom}_{\mathcal{C}}(A,B)$ so that we really have composition of functions here. 
        \item These functions are also associative such that for $f \in \text{Hom}_{\mathcal{C}}(A,B), g \in \text{Hom}_{\mathcal{C}}(B,C), h \in \text{Hom}_{\mathcal{C}}(C,D)$ we have $h \circ (g \circ f) = (h \circ g) \circ f$ and the same for units $\text{id}_B \circ f = f \circ \text{id}_A = f$.
    \end{enumerate}
\end{definition}

\noindent
So what's cool about this is that this is basically a group without the inverses requirement. It becomes hard to rigorously define what this collection of objects really is, and especially to decide if it is a set, so they handwaved this, but we'll come back to it later here. Now, we continue with an example.

\begin{example}
    Placeholder
\end{example}

\end{document}
